% The nabia Python package on the wheel: the library from Python, backtests, market data, analytics.

\chapter{Portfolio backtests: \texorpdfstring{\code{nabia.backtest}}{nabia.backtest}}\label{p:sec:bindings}
\code{nabia.portfolio} runs the C++ portfolio engine through \code{\_rb\_backtest}, a module of
the wheel's package \code{nabia\_python\_bridge}. The module reuses the command-line
application's JSON parser: Python and the command line accept the same parameters with the same
error messages (Chapter~\ref{c:sec:rbbindings}). Data crosses the boundary
in simple forms: dates as \code{int64} days since 1970-01-01, prices as a \code{float64} matrix,
and parameters as JSON text.

Two functions are exposed:
\begin{itemize}
  \item \code{nabia.backtest(prices, params, budget\_schedule=None)} returns a result with
        pandas frames and metrics;
  \item \code{nabia.latest\_weights(prices, params)} solves only the last rebalance: the target weights
        to hold now.
\end{itemize}
\snippet{py_backtest_inputs}{Python}{pandas to C++}
\begin{codepointers}
  \codeitem{py_backtest}{Python \code{backtest}}
  \codeitem{py_latest_weights}{Python \code{latest\_weights}}
  \codeitem{py_to_days}{dates to day numbers}
\end{codepointers}

%==============================================================================


\chapter{Market data}\label{p:sec:marketdata}

\section{Universe and sources}
The default download universe has four asset classes:
\begin{itemize}
  \item equity indices: S\&P 500, Nasdaq 100, Russell 2000, Euro Stoxx 50, DAX, FTSE 100,
        Nikkei 225, Hang Seng;
  \item US Treasury futures: 2Y, 5Y, 10Y, Bond, Ultra Bond;
  \item commodity futures: WTI, Brent, natural gas, gold, silver, copper, corn, wheat,
        soybeans;
  \item crypto against USD: BTC, ETH, SOL, XRP, ADA, LTC.
\end{itemize}
The AI portfolio uses the nine assets of Table~\ref{p:tab:universe-default}.

\begin{table}[ht]\centering\small
\caption{A default nine-asset portfolio universe.}\label{p:tab:universe-default}
\begin{tabular}{@{}lll@{}}
\toprule
Symbol & Asset & Class \\
\midrule
\code{\textasciicircum GSPC} & S\&P 500 & equity index \\
\code{\textasciicircum STOXX50E} & Euro Stoxx 50 & equity index \\
\code{\textasciicircum N225} & Nikkei 225 & equity index \\
\code{ZN=F} & US 10-Year T-Note future & bond future \\
\code{ZB=F} & US Treasury Bond future & bond future \\
\code{BZ=F} & Brent crude oil future & commodity \\
\code{GC=F} & Gold future & commodity \\
\code{HG=F} & Copper future & commodity \\
\code{BTC/USD} & Bitcoin & crypto \\
\bottomrule
\end{tabular}
\end{table}

\paragraph{Sources.}
\begin{itemize}
  \item \textbf{Yahoo Finance} (\code{yfinance}) provides indices and futures, in batches.
        Futures are front-month continuous series, which jump at contract rolls. WTI
        settled negative on 2020-04-20, where simple returns are undefined, so the
        portfolio uses Brent.
  \item \textbf{Crypto exchanges} (\code{ccxt}, Coinbase by default) provide daily
        candles from public data, with no key. History is paged in windows of 300 candles.
        Windows before a coin's listing date are empty and are stepped over. Paging stops
        if the exchange ignores the start date.
\end{itemize}
\snippet{crypto_fetch}{Python}{Crypto daily candles with pagination}

\section{Storage and incremental updates}
Prices are stored per (symbol, date): open, high, low, close, adjusted close and volume.
\begin{itemize}
  \item \textbf{Upserts.} Writes use \code{INSERT ... ON CONFLICT DO UPDATE}, in the SQLite
        or PostgreSQL dialect, so re-downloading is idempotent and picks up revisions.
  \item \textbf{Incremental updates.} A stored symbol is re-fetched from five days before
        its last date, which corrects recent revisions.
  \item \textbf{Batching.} Symbols are grouped by source and start date.
\end{itemize}
\snippet{upsert_rows}{Python}{Portable upsert}
\snippet{update_market_data}{Python}{Incremental download}
\begin{codepointers}
  \codeitem{universe}{universe}
  \codeitem{yahoo_fetch}{Yahoo source}
  \codeitem{load_prices}{loading a price matrix}
\end{codepointers}

\chapter{Analytics on price frames}\label{p:sec:stats}
\code{nabia.analytics} computes descriptive statistics and trend signals on any
pandas frame of prices (dates by assets).

All statistics use daily close prices up to \code{as\_of}, over
\code{stats\_history\_days} (400) calendar days. Volatilities are annualized with
$\sqrt{252}$, standard deviations are sample standard deviations, and returns are simple
returns.
\begin{description}
  \item[Aligned returns] keep the dates where all assets have a price, then take returns.
        They are used for correlations only. Single-asset statistics use each asset's own
        history.
  \item[Correlation] is the Pearson correlation matrix $\rho$ of the last $w$ aligned
        returns (60 by default, at least 3). The \emph{average correlation} is
        \[
          \bar\rho=\frac{1}{n(n-1)}\sum_{i\ne j}\rho_{ij}.
        \]
  \item[Rolling volatility] is
        $\sigma^{(w)}=\mathrm{sd}\big(r_{T-w+1..T}\big)\sqrt{252}$ for
        $w\in\{20,60,120\}$. The \emph{volatility percentile} is computed from the series
        $v_s$ of rolling 20-day standard deviations over the last 252 dates:
        \[
          \pi=\frac{1}{|v|}\sum_s\mathbf 1[v_s\le v_T].
        \]
        So $\pi=1$ means the highest of the year.
  \item[Past returns] $P_T/P_{T-n}-1$ for $n\in\{21,63,126,252\}$ observations
        (1, 3, 6 and 12 months).
  \item[Moving-average distance] $P_T/\overline P_n-1$, where $\overline P_n$ is the mean
        of the last $n\in\{50,200\}$ prices.
  \item[Volatility-scaled 12-1 momentum]
        \[
          m=\frac{P_{T-21}/P_{T-252}-1}{\mathrm{sd}(r_{T-251..T})\sqrt{252}} .
        \]
        It is the 11-month return ending one month ago, which skips the short-term
        reversal, measured in units of the asset's own annual volatility.
  \item[Drawdowns] over the last 253 prices: $D_t=1-P_t/\max_{s\le t}P_s$. The tools
        report the current value and the maximum.
  \item[Market regime] combines five numbers:
        \begin{itemize}
          \item $\bar\rho$ over 60 and over 252 aligned returns;
          \item the median volatility percentile;
          \item the share of assets above their 200-day average;
          \item the share of assets with a positive 12-month return.
        \end{itemize}
\end{description}
\snippet{trend_signals}{Python}{Trend signals}
\snippet{rolling_volatility}{Python}{Volatilities and their percentile}
\snippet{market_regime}{Python}{Market regime}
\begin{codepointers}
  \codeitem{correlation_matrix}{correlation matrix}
  \codeitem{average_correlation}{average correlation}
\end{codepointers}
